% Learning InitechBase, Chapter 4: Finding Records (AN-16).  Every panel
% replayed (examples/l4-*.dot); sources cmd/query.c (LOCATE/CONTINUE/
% SEEK/FIND), samir_main.c repl_do_use (INDEX), fs/ndx.c, cmd/set.c.
\chapter{Finding Records}
\Gates{test-pubs-examples, test-samir-query, test-ndx-seek}
Initech's telephone list is no use if you must read every record to find one
person.  In this lesson InitechBase finds the records for you: first by
reading the database until a condition is met, then at once, through an
index.

\begin{ChapterCovers}
\item Searching with \cmd{LOCATE} and \cmd{CONTINUE}
\item \fn{FOUND()}
\item Index files
\item Finding a key with \cmd{SEEK} and \cmd{FIND}
\item SET EXACT
\item Choosing the order with SET ORDER
\end{ChapterCovers}

\ATEHead{Preparing for This Chapter}
Start InitechBase with a fresh copy of the sample files.  Each part of the
lesson begins with \cmd{USE}, so you need not open a database yet.

\ATEHead{Searching with LOCATE}
\cmd{LOCATE} reads the records from the first and stops at the first record
for which the condition after \kw{FOR} is true.  \fn{FOUND()} tells you
whether a record was found.  \cmd{CONTINUE} goes on to the next.
\begin{Steps}
\item Type \type{USE STAFF}.
\item Type \type{LOCATE FOR DEPT = 'ACCT'}, then \type{? RECNO(), LASTNAME, FOUND()}.
\item Type \type{CONTINUE}, then \type{? RECNO(), LASTNAME, FOUND()}.
\item Type \type{CONTINUE}, then \type{? FOUND(), EOF()}.
\ATEResponse{There is no third record: the pointer is at the end of the file.}
\end{Steps}
\DotExample{examples/l4-locate}
After \cmd{LOCATE}, the record found is the current record, so \cmd{DISPLAY}
shows it:
\DotExample{examples/l4-locate-2}

\ATEHead{Index Files}
\cmd{LOCATE} reads every record until it finds one, which takes time in a large
database.  An index file keeps the records in the order of a key and lets
InitechBase go straight to a key.  The sample diskette holds four index files
for the staff database: \file{STAFFNAM.NDX} on the last name,
\file{STAFFNO.NDX} on the employee number, \file{STAFFDPT.NDX} on the
department and \file{STAFFSAL.NDX} on the salary.

An index file is opened with its database, by \kw{INDEX} in \cmd{USE}.  The
records are then taken in the order of the index.
\begin{Steps}[resume]
\item Type \type{USE STAFF INDEX STAFFNAM}.
\item Type \type{LIST LASTNAME, FIRSTNAME, DEPT}.
\ATEResponse{The employees are listed in alphabetical order of their last names.}
\end{Steps}
\DotExample{examples/l4-index}
The record numbers have not changed: the index changes only the order in which
InitechBase takes the records.

\ATEHead{Finding a Key}
\cmd{SEEK} looks up a value in the index and moves to the first record that has
it.
\begin{Steps}[resume]
\item Type \type{SEEK 'OKAFOR'}, then \type{? FOUND(), RECNO()}.
\item Type \type{DISPLAY LASTNAME, FIRSTNAME, EXT}.
\end{Steps}
\DotExample{examples/l4-seek}
When no key matches, \fn{FOUND()} is false and the pointer is at the end of the
file:
\DotExample{examples/l4-notfound}
\cmd{FIND} does the same with the value typed without quotation marks:
\DotExample{examples/l4-find}

\ATEHead{Part of a Key}
\cmd{SEEK} finds a key that begins with the value you give, so you need type
only the first letters of a long name.  SET EXACT ON makes it find only a key
that is equal to the value.
\begin{Steps}[resume]
\item Type \type{SEEK 'NAG'}, then \type{? LASTNAME}.
\item Type \type{SET EXACT ON}, then \type{SEEK 'NAG'}, then \type{? FOUND()}.
\end{Steps}
\DotExample{examples/l4-exact}
Type \type{SET EXACT OFF} before you go on.

\ATEHead{Several Index Files}
You can open several index files with a database.  The first named is the
master index: it gives the order and is the index \cmd{SEEK} uses.  The others
are kept up to date as you change records.  SET ORDER chooses another master
index.
\begin{Steps}[resume]
\item Type \type{USE STAFF INDEX STAFFNAM, STAFFDPT}.
\item Type \type{SET ORDER TO 2}, then \type{GO TOP}.
\item Type \type{LIST NEXT 5 LASTNAME, DEPT}.
\end{Steps}
\DotExample{examples/l4-order}
With the salary index, the last record is the highest salary:
\DotExample{examples/l4-salary}
\begin{Tip}
After SET ORDER, type \type{GO TOP} so that the current record is the first in
the new order.
\end{Tip}

\ATEHead{Where Do You Go From Here?}
You can find any record by a condition or by its key.  In Chapter~5 you change
the records you find.

\ATEHead{Summary}
\begin{enumerate}
\item \cmd{LOCATE FOR} \xp{condition} finds the first record that meets the
  condition; \cmd{CONTINUE} finds the next; \fn{FOUND()} reports success.
\item \cmd{USE} \xp{database} \kw{INDEX} \xp{index list} opens index files; the
  first is the master index.
\item \cmd{SEEK} \xp{expression} and \cmd{FIND} \xp{characters} find a key in the
  master index.
\item SET EXACT ON makes \cmd{SEEK} find only an equal key.
\item SET ORDER TO \xp{n} makes the \xp{n}th index file the master index.
\end{enumerate}
